% Part: lambda-calculus
% Chapter: introduction
% Section: minimization

\documentclass[../../../include/open-logic-section]{subfiles}

\begin{document}

\olfileid[hi]{lam}{int}{min}
\olsection{\usetoken{S}{lambda definable} फलन न्यूनीकरण के अधीन संवृत हैं}

\begin{lem}
मान लें $f(x,y)$ !!{lambda definable} है। $g$ को इस प्रकार परिभाषित करें:
\[
g(x) \simeq \umin{y}{f(x,y)}.
\]
तब $g$ !!{lambda definable} है।
\end{lem}

\begin{proof}
मोटे तौर पर विचार यह है। $x$ दिए होने पर हम अचल-बिंदु लैम्ब्डा पद $Y$ का
उपयोग ऐसा फलन $h_x(n)$ परिभाषित करने के लिए करेंगे जो~$n$ से आरंभ करके
किसी~$y$ को खोजता है; तब $g(x)$ केवल $h_x(0)$ है। फलन~$h_x$ को किसी
अचल-बिंदु समीकरण के हल के रूप में व्यक्त किया जा सकता है:
\[
h_x(n) \simeq
\begin{cases}
n & \text{यदि $f(x,n) = 0$} \\
h_x(n+1) & \text{अन्यथा।}
\end{cases}
\]

अब विवरण देखें। चूँकि $f$ आदिम पुनरावर्ती है, वह किसी पद $F$ द्वारा
!!{lambda defined} है। स्मरण करें कि हमारे पास ऐसा लैम्ब्डा पद~$D$ भी है
कि $D(M, N, \bar{0}) \red M$ और $D(M, N, \bar{1}) \red N$। अभी $x$ को
नियत रखते हुए, $h_x$ को !!{lambda define} करने के लिए हम~$x$ पर निर्भर ऐसा
पद~$H$ खोजना चाहते हैं जो
\[
H(\num{n}) \equiv D(\num{n}, H(S(\num{n})), F(x, \num{n})).
\]
अचल-बिंदु पद~$Y$ का उपयोग करके हम ऐसा कर सकते हैं। पहले $U$ यह पद हो:
\[
\lambd[h][\lambd[z][D(z,(h(Sz)),F(x,z))]],
\]
और फिर $H$ पद~$YU$ हो। ध्यान दें कि~$H$ में एकमात्र मुक्त चर~$x$ है। हम
दिखाएँ कि $H$ ऊपर के समीकरण को संतुष्ट करता है।

$Y$ की परिभाषा से,
\[
H = YU \equiv U(YU) = U(H).
\]
विशेषतः, प्रत्येक प्राकृतिक संख्या~$n$ के लिए,
\begin{align*}
H(\num{n}) & \equiv U(H, \num{n}) \\
& \red D(\num{n}, H(S(\num{n})), F(x, \num{n})),
\end{align*}
जैसा अपेक्षित था। ध्यान दें कि अंतिम पंक्ति में~$x$ के स्थान पर कोई अंक
$\num{m}$ प्रतिस्थापित करने पर, यदि $F(\num{m}, \num{n})$ अपचयित होकर
$\num{0}$ बनता है तो व्यंजक अपचयित होकर $\num{n}$ बनता है; और यदि
$F(\num{m}, \num{n})$ अपचयित होकर कोई अन्य अंक बनता है, तो वह अपचयित होकर
$H(S(\num{n}))$ बनता है।

प्रमाण पूरा करने के लिए $G$, ~$\lambd[x][H(\num 0)]$ हो। तब $G$, ~$g$ को
!!{lambda define} करता है; दूसरे शब्दों में, प्रत्येक~$m$ के लिए यदि
$g(m)$ परिभाषित है तो $G(\num m)$ अपचयित होकर~$\overline {g(m)}$ बनता है,
और अन्यथा उसका कोई प्रसामान्य रूप नहीं होता।
\end{proof}

\end{document}
